\scanpage{197}
penser que pour les \unsure{pts différents} à val. dans $\unsure{k_0}\subset k$, les deux foncteurs obtenus, soient \unsure{isom.} (pas plus que dans le cas (ii)). Il n'y a plus en général d'opération de $\pi_1(X)$ sur $\xi_{x,\infty}$, mais encore comme dans (i) les autom. de $k$ sur $\unsure{k(x)}$, opèrent dans $\xi_{x,\infty}$ \unsure{$\frac12$} \uline{linéaires} (donnée de descente). Un autom. de $k$ qui n'est pas l'identité sur $k$ définit $\xi_{x,\infty}^{\unsure{\sigma}}$ qui à priori n'est plus isomorphe à $\xi_{x,\infty}$.

Un cas intéressant est celui où $k=\mathbf Q$, on trouve alors un \deleted{point} $\xi_{x,\infty}\in T_{\unsure{\alpha}}(\mathbf Q)$, i.e. une trivialisation de $T_{\mathbf Q}$. Deux trivialisations \unsure{à val. dans} $T_{\mathbf Q}$ ne sont en général pas \unsure{isom.} à priori du moins. De plus, il n'y a pas d'isom. \uline{canonique} privilégié entre
\[
\xi_{x,\infty}\quad\unsure{\text{de}}\quad\xi_{x,\infty}\otimes_{\mathbf Q}\mathbf Q_\ell
\quad\text{et}\quad\xi_{\bar x,\ell}\otimes_{\mathbf Z_\ell}\mathbf Q_\ell\quad-
\]
\deleted{\unsure{bien que pas}} et comme le premier à priori ne provient pas d'un objet de $\unsure{T_{\alpha}}(\mathbf Z_\ell)$ mais seulement de $T(\mathbf Q_\ell)$, il n'y a pas lieu de penser qu'il y ait \unsure{même} un isom. Ainsi, tt pt $x/\mathbf Q$ définit une classe de \unsure{ch} dans les $H^{\unsure{2|1}}(\mathbf Q_\ell,G_\ell^{\circ(\unsure{\alpha})})$. Dans le cas général, il faut remplacer $\mathbf Q_\ell$

\margin{en fait, ils sont déjà \unsure{isom.} sur les motifs de poids et niveau $0$}
\margin{196}
